\documentclass[11pt, letterpaper]{article}
\usepackage{afp-style}
\afpdoctype{Activation Guide}

\begin{document}

\afptitlepage{Chapter Activation\\Guide}{A Step-by-Step Playbook for Starting Your Chapter}

\tableofcontents
\newpage

\section{Before You Begin}

Starting a chapter of the Anti-Federalist Party is intentionally simple. We practice the decentralization we preach --- you do not need permission from a national headquarters. You need five committed people and the willingness to do the work.

\afpcallout{This guide assumes you are starting from zero: no existing political organization, no funding, no infrastructure. By the end of this process, you will have a functioning, chartered chapter capable of endorsing candidates, passing resolutions, and participating in the national movement.}

\subsection{What You Need}
\begin{itemize}
  \item A minimum of \textbf{five (5) founding members} who live in your geographic area.
  \item A copy of the \textbf{Party Constitution} (available at anti-federalists.com/documents).
  \item A copy of the \textbf{Chapter Charter} template.
  \item A meeting space (a living room, library meeting room, coffee shop, or virtual platform).
  \item Approximately 2--3 hours for your organizational meeting.
\end{itemize}


\section{Phase 1: Gather Your Founding Members}

\subsection{Finding People}
\begin{enumerate}
  \item \textbf{Start with who you know.} The first five members will likely come from your existing network --- friends, neighbors, coworkers, or fellow community members who share Anti-Federalist values.
  \item \textbf{Use local channels.} Post on community bulletin boards (physical and digital), neighborhood apps like Nextdoor, or local social media groups.
  \item \textbf{Attend existing meetings.} City council meetings, county commission meetings, and school board meetings are where you'll find people who care about local governance.
  \item \textbf{Host an interest meeting.} Before your formal organizational meeting, host a casual gathering to discuss the Party's principles and gauge interest.
\end{enumerate}

\subsection{The Pitch}
When recruiting founding members, focus on three points:
\begin{enumerate}
  \item \textbf{Local power.} ``We believe the people in this community should make the decisions that affect this community --- not politicians in Washington or the state capital.''
  \item \textbf{Non-partisan.} ``We're not left or right. We're local. We critique both parties equally because both parties centralize power.''
  \item \textbf{Action-oriented.} ``We're not a discussion group. We endorse candidates, pass resolutions, and participate in local governance.''
\end{enumerate}


\section{Phase 2: The Organizational Meeting}

\subsection{Agenda}
Your first official meeting should follow this agenda:

\begin{enumerate}
  \item \textbf{Welcome and introductions} (10 minutes)
  \item \textbf{Review of Party Principles} --- Read Article~II of the Constitution aloud. Discuss. (20 minutes)
  \item \textbf{Adoption of the Constitution} --- Motion to adopt the Party Constitution as the chapter's governing document. Vote. (5 minutes)
  \item \textbf{Election of officers} --- Nominate and elect: Chair, Vice Chair, Secretary, Treasurer. (30 minutes)
  \item \textbf{Completion of the Chapter Charter} --- Fill in the Charter template. All founding members sign. (15 minutes)
  \item \textbf{Operational decisions:}
    \begin{itemize}
      \item Set regular meeting schedule (day, time, frequency).
      \item Choose communication platform (Signal, email list, etc.).
      \item Decide on dues (if any).
    \end{itemize}
    (20 minutes)
  \item \textbf{First business} --- Identify one concrete action item for the chapter's first 30 days. (15 minutes)
  \item \textbf{Adjournment}
\end{enumerate}


\section{Phase 3: Formalize Your Chapter}

\subsection{Submit Your Charter}
Email your signed Chapter Charter to \texttt{chapters@anti-federalists.com} with:
\begin{itemize}
  \item Scanned or photographed copy of the signed Charter.
  \item Contact email for the Chair and Secretary.
  \item Brief description of your geographic jurisdiction.
\end{itemize}

\subsection{Establish Your Finances}
\begin{enumerate}
  \item Open a bank account in the chapter's name at a \textbf{local credit union} or community bank (practice what we preach).
  \item Require two signatures on all expenditures over \$100.
  \item The Treasurer maintains a ledger of all income and expenses.
  \item Financial records are open to any member at any time.
\end{enumerate}

\subsection{Set Up Communications}
\begin{enumerate}
  \item \textbf{Encrypted messaging:} Signal group for leadership communication.
  \item \textbf{Email list:} For meeting notices and member communication.
  \item \textbf{Optional:} Social media presence on platforms of your choice.
  \item \textbf{Avoid:} Centralized platforms that can deplatform you. Own your communication infrastructure where possible.
\end{enumerate}


\section{Phase 4: First 90 Days}

\subsection{Month 1: Establish Presence}
\begin{itemize}
  \item Hold your first regular meeting.
  \item Attend a local government meeting (city council, county commission, or school board) as a group.
  \item Identify one local issue your chapter cares about.
  \item Begin recruiting beyond your founding members.
\end{itemize}

\subsection{Month 2: Engage}
\begin{itemize}
  \item Draft your first resolution on a local issue.
  \item Invite a local elected official to speak at a chapter meeting.
  \item Establish a relationship with local media (community newspaper, local radio).
  \item Host a public event --- a community forum, candidate Q\&A, or issue discussion.
\end{itemize}

\subsection{Month 3: Act}
\begin{itemize}
  \item Pass your first resolution and submit it to relevant local authorities.
  \item Begin identifying potential candidates for local office who align with Party principles.
  \item Connect with other Anti-Federalist chapters in your state.
  \item Submit your first quarterly report to the national registry.
\end{itemize}


\section{Phase 5: Ongoing Operations}

\subsection{Meeting Best Practices}
\begin{itemize}
  \item Always have a written agenda distributed in advance.
  \item Start and end on time.
  \item Use Robert's Rules of Order (simplified) for motions and votes.
  \item Record minutes and distribute within 14 days.
  \item Rotate facilitation duties to develop leadership across the chapter.
\end{itemize}

\subsection{Membership Growth}
\begin{itemize}
  \item Every member recruits one new member per quarter (the ``Each One, Reach One'' standard).
  \item Host quarterly public events open to non-members.
  \item Maintain a welcoming, non-ideological-purity-test culture. Agreement on Principles is sufficient.
\end{itemize}

\subsection{Candidate Development}
\begin{itemize}
  \item Identify local offices with upcoming elections (school board, city council, county commission, water board, etc.).
  \item Recruit and support candidates from within the chapter.
  \item Focus on winnable races. A school board seat won is worth more than a Congressional campaign lost.
\end{itemize}


\section{Resources}

\begin{itemize}
  \item \textbf{Party Constitution} --- The governing document of the Party.
  \item \textbf{Chapter Charter Template} --- Fill-in charter for chapter formation.
  \item \textbf{Membership Application} --- Standard form for new members.
  \item \textbf{Meeting Minutes Template} --- Standardized meeting record format.
  \item \textbf{Resolution Template} --- Formal resolution format for chapter positions.
  \item \textbf{Treasurer's Report Template} --- Financial reporting format.
\end{itemize}

\vspace{1em}
All documents available at: \texttt{anti-federalists.com/documents}

\vspace{2cm}

\afpcertification{Chapter Activation Guide}

\vspace{1cm}

\begin{center}
  {\itshape\color{afp-gray}\sffamily
    The revolution starts in your neighborhood.\\
    It always has.
  }
\end{center}

\end{document}
